% !tex ts-program = pdflatex
% !tex encoding = utf-8 unicode

% this is a simple template for a latex document using the "article" class.
% see "book", "report", "letter" for other types of document.

\documentclass[11pt]{article} % use larger type; default would be 10pt
\usepackage[
    type={CC},
    modifier={by-nc-sa},
    version={3.0},
]{doclicense}

\usepackage[utf8]{inputenc} % set input encoding (not needed with xelatex)
\usepackage[strict]{changepage}
    
    %%% examples of article customizations
    % these packages are optional, depending whether you want the features they provide.
    % see the latex companion or other references for full information.
    
    
    %%% page dimensions
\usepackage{geometry} % to change the page dimensions
\geometry{a4paper} % or letterpaper (us) or a5paper or....
    % \geometry{margin=2in} % for example, change the margins to 2 inches all round
    % \geometry{landscape} % set up the page for landscape
    %   read geometry.pdf for detailed page layout information
 \usepackage{numdef}
   
\usepackage{graphicx} % support the \includegraphics command and options
    % some of the article customisations are relevant for this class
\usepackage{amsmath,amsthm}
\usepackage{amsfonts} % math fonts such as \mathbb{}
\usepackage{amssymb} % \therefore
% \usepackage{bickham}
\usepackage{hyperref}
\usepackage{cryptocode}
\usepackage{framed} 
    % \usepackage[parfill]{parskip} % activate to begin paragraphs with an empty line rather than an indent
    
    %%% packages
\usepackage{booktabs} % for much better looking tables
\usepackage{array} % for better arrays (eg matrices) in maths
\usepackage{paralist} % very flexible & customisable lists (eg. Enumerate/itemize, etc.)
\usepackage{verbatim} % adds environment for commenting out blocks of text & for better verbatim
\usepackage{subfig} % make it possible to include more than one captioned figure/table in a single float
    % % these packages are all incorporated in the memoir class to one degree or another...
\usepackage{mathrsfs}
\usepackage{booktabs}
\usepackage{makecell}
\usepackage{adjustbox}

\usepackage{pgfplots}

\renewcommand\theadalign{bc}
\renewcommand\theadfont{\bfseries}
\renewcommand\theadgape{\Gape[4pt]}
  
    %%% headers & footers
\usepackage{fancyhdr} % this should be set after setting up the page geometry
\pagestyle{fancy} % options: empty , plain , fancy
\renewcommand{\headrulewidth}{0pt} % customise the layout...
\lhead{}\chead{}\rhead{}
\lfoot{}\cfoot{\thepage}\rfoot{}
    
    
    %%% section title appearance
\usepackage{sectsty}
\allsectionsfont{\sffamily\mdseries\upshape} % (see the fntguide.pdf for font help)
    % (this matches context defaults)
    
    %%% toc (table of contents) appearance
\usepackage[nottoc,notlof,notlot]{tocbibind} % put the bibliography in the toc
\usepackage[titles,subfigure]{tocloft} % alter the style of the table of contents
\renewcommand{\cftsecfont}{\rmfamily\mdseries\upshape}
\renewcommand{\cftsecpagefont}{\rmfamily\mdseries\upshape} % no bold!

    %%% end article customizations
    \newcommand{\code}[1]{\texttt{#1}}
\newcommand\tstrut{\rule{0pt}{2.6ex}}         % = `top' strut
\newcommand\bstrut{\rule[-0.9ex]{0pt}{0pt}}   % = `bottom' strut
\newcommand{\bgamma}{\boldsymbol{\gamma}}
\newcommand{\bsigma}{\boldsymbol{\sigma}}
\newcommand{\plonk}{\ensuremath{\mathcal{P} \mathfrak{lon}\mathcal{K}}\xspace}

% \newcommand{\plonk}{\ensuremath{\mathtt{PLONK}}\xspace}
%%% the "real" document content comes below...
\newcommand{\papertitle}{Ironwood's knowledge soundness proof of IPA in the Algebraic Group Model}
%\newcommand{\authorname}}
\newcommand{\company}{}
\title{ \bf \papertitle \\[0.72cm]}
\author{ Ariel Gabizon}


% 	\large{\authorname} \\[0.5cm] \large{\company}
% 	\\ {DRAFT}
%}
    %\date{} % activate to display a given date or no date (if empty),
  
           % otherwise the current date is printed 
	
\ProvidesPackage{numdef}

%% Ariel Macros:
% \num\newcommand{\G1}{\ensuremath{{\mathbb G}_1}\xspace}
\newcommand{\Gi}{\ensuremath{{\mathbb G}_i}\xspace}
\newcommand{\G}{\ensuremath{{\mathbb G}}\xspace}
\newcommand{\Gstar}{\ensuremath{{\mathbb G}^*}\xspace}
\newcommand{\x}{\ensuremath{\mathbf{x}}\xspace}
\newcommand{\z}{\ensuremath{\mathbf{z}}\xspace}
\newcommand{\X}{\ensuremath{\mathbf{X}}\xspace}

% \num\newcommand{\G2}{\ensuremath{{\mathbb G}_2}\xspace}
%\num\newcommand{\G11}{\ensuremath{\G1\setminus \set{0} }\xspace}
%\num\newcommand{\G21}{\ensuremath{\G2\setminus \set{0} }\xspace}
\newcommand{\grouppair}{\ensuremath{G^*}\xspace}
\newcommand{\prvperm}{\ensuremath{\mathrm{P_{\mathsf{\sigma}}}}\xspace}
\newcommand{\verperm}{\ensuremath{\mathrm{V_{\mathsf{\sigma}}}}\xspace}

\newcommand{\Gt}{\ensuremath{{\mathbb G}_t}\xspace}
\newcommand{\F}{\ensuremath{\mathbb F}\xspace}
\newcommand{\Fstar}{\ensuremath{\mathbb F^*}\xspace}

\newcommand{\help}[1]{$#1$-helper\xspace}
\newcommand{\randompair}[1]{\ensuremath{\mathsf{randomPair}(#1)}\xspace}
\newcommand{\pair}[1]{$#1$-pair\xspace}
\newcommand{\pairs}[1]{$#1$-pairs\xspace}
\newcommand{\chalpoint}{\ensuremath{\mathfrak{z}}\xspace}

\newcommand{\pairone}[1]{\G1-$#1$-pair\xspace}
\newcommand{\pairtwo}[1]{\G2-$#1$-pair\xspace}
\newcommand{\sameratio}[2]{\ensuremath{\mathsf{SameRatio}(#1,#2)}\xspace}
\newcommand{\vecc}[2]{\ensuremath{(#1)_{#2}}\xspace}
\newcommand{\players}{\ensuremath{[n]}\xspace}
\newcommand{\adv}{\ensuremath{\mathcal A}\xspace}
\newcommand{\advprime}{\ensuremath{{\mathcal A}'}\xspace}
\newcommand{\extprime}{\ensuremath{E'}\xspace}
\newcommand{\advrand}{\ensuremath{\mathsf{rand}_{\adv}}\xspace}
% \num\newcommand{\srs1}{\ensuremath{\mathsf{srs_1}}\xspace}
% \num\newcommand{\srs2}{\ensuremath{\mathsf{srs_2}}\xspace}
\newcommand{\srs}{\ensuremath{\mathsf{srs}}\xspace}
\newcommand{\srsi}{\ensuremath{\mathsf{srs_i}}\xspace}
\newcommand{\com}{\ensuremath{\mathsf{com}}\xspace}
\newcommand{\comperm}{\ensuremath{\mathsf{com_{\sigma}}}\xspace}
\newcommand{\cm}{\ensuremath{\mathsf{cm}}\xspace}
\newcommand{\cmsig}{\ensuremath{\mathsf{cm_\sigma}}\xspace}
\newcommand{\open}{\ensuremath{\mathsf{open}}\xspace}
\newcommand{\openperm}{\ensuremath{\mathsf{open_{\sigma}}}\xspace}
\newcommand{\sigof}[1]{\ensuremath{\sigma(#1)}\xspace}
\newcommand{\proverexp}{\ensuremath{\mathsf{e}}\xspace}
\newcommand{\reducedelems}{\ensuremath{\mathsf{r}}\xspace}

\newcommand{\ci}{\ensuremath{\mathrm{CI}}\xspace}
\renewcommand{\deg}{\ensuremath{\mathrm{deg}}\xspace}
\newcommand{\pairvec}[1]{$#1$-vector\xspace}
\newcommand{\Fq}{\ensuremath{\mathbb{F}_q}\xspace}
\newcommand{\randpair}[1]{\ensuremath{\mathsf{rp}_{#1}}\xspace}
\newcommand{\randpairone}[1]{\ensuremath{\mathsf{rp}_{#1}^{1}}\xspace}
\newcommand{\abase}{\ensuremath{A_{\mathrm{\mathbf{0}}}}\xspace}
\newcommand{\bbase}{\ensuremath{B_{\mathrm{\mathbf{0}}}}\xspace}
\newcommand{\cbase}{\ensuremath{C_{\mathrm{\mathbf{0}}}}\xspace}

\newcommand{\amid}{\ensuremath{A_{\mathrm{mid}}}\xspace}
\newcommand{\bmid}{\ensuremath{B_{\mathrm{mid}}}\xspace}
\newcommand{\cmid}{\ensuremath{C_{\mathrm{mid}}}\xspace}

\newcommand{\negl}{\ensuremath{\mathsf{negl}(\lambda)}\xspace}
\newcommand{\randpairtwo}[1]{\ensuremath{\mathsf{rp_{#1}^2}}\xspace}%the randpair in G2
\newcommand{\nilp}{\ensuremath{\mathscr N}\xspace}
\newcommand{\groupgen}{\ensuremath{\mathscr G}\xspace}
\newcommand{\qap}{\ensuremath{\mathscr Q}\xspace}
\newcommand{\polprot}[4]{$(#1,#2,#3,#4)$-polynomial protocol}
\newcommand{\rangedprot}[5]{$#5$-ranged $(#1,#2,#3,#4)$-polynomial protocol}

\newcommand{\rej}{\ensuremath{\mathsf{rej}}\xspace}
\newcommand{\acc}{\ensuremath{\mathsf{acc}}\xspace}
\newcommand{\res}{\ensuremath{\mathsf{res}}\xspace}
\newcommand{\sha}[1]{\ensuremath{\mathsf{COMMIT}(#1)}\xspace}
 \newcommand{\shaa}{\ensuremath{\mathsf{COMMIT}}\xspace}
 \newcommand{\comm}[1]{\ensuremath{\mathsf{comm}_{#1}}\xspace}
 \newcommand{\defeq}{:=}

\newcommand{\B}{\ensuremath{\vec{B}}\xspace}
\newcommand{\dom}{\ensuremath{H}\xspace}
\newcommand{\C}{\ensuremath{\vec{C}}\xspace}
\newcommand{\Btwo}{\ensuremath{\vec{B_2}}\xspace}
\newcommand{\treevecsimp}{\ensuremath{(\tau,\rho_A,\rho_A \rho_B,\rho_A\alpha_A,\rho_A\rho_B\alpha_B, \rho_A\rho_B\alpha_C,\beta,\beta\gamma)}\xspace}% The sets of elements used in simplifed relation tree in main text body
\newcommand{\rcptc}{random-coefficient subprotocol\xspace}
\newcommand{\rcptcparams}[2]{\ensuremath{\mathrm{RCPC}(#1,#2)}\xspace}
\newcommand{\verifyrcptcparams}[2]{\ensuremath{\mathrm{\mathsf{verify}RCPC}(#1,#2)}\xspace}
\newcommand{\randadv}{\ensuremath{\mathsf{rand}_{\adv}}\xspace}
 \num\newcommand{\ex1}[1]{\ensuremath{ #1\cdot g_1}\xspace}
 \num\newcommand{\ex2}[1]{\ensuremath{#1\cdot g_2}\xspace}
 \newcommand{\pr}{\mathrm{Pr}}
 \newcommand{\powervec}[2]{\ensuremath{(1,#1,#1^{2},\ldots,#1^{#2})}\xspace}
 \newcommand{\partition}{\ensuremath{{\mathcal T}}\xspace}
 \newcommand{\partof}[1]{\ensuremath{{\partition_{#1}}}\xspace}
\num\newcommand{\out1}[1]{\ensuremath{\ex1{\powervec{#1}{d}}}\xspace}
\num\newcommand{\out2}[1]{\ensuremath{\ex2{\powervec{#1}{d}}}\xspace}
 \newcommand{\nizk}[2]{\ensuremath{\mathrm{NIZK}(#1,#2)}\xspace}% #2 is the hash concatenation input
 \newcommand{\verifynizk}[3]{\ensuremath{\mathrm{VERIFY\mhyphen NIZK}(#1,#2,#3)}\xspace}
\newcommand{\protver}{protocol verifier\xspace} 
\newcommand{\hash}{\ensuremath{\mathcal{H}}\xspace} 
\newcommand{\mulgroup}{\ensuremath{\F^*}\xspace}
\newcommand{\lag}[1]{\ensuremath{L_{#1}}\xspace} 
\newcommand{\sett}[2]{\ensuremath{\set{#1}_{#2}}\xspace}
\newcommand{\omegaprod}{\ensuremath{\alpha_{\omega}}\xspace}
\newcommand{\lagvec}[1]{\ensuremath{\mathrm{LAG}_{#1}}\xspace}
\newcommand{\trapdoor}{\ensuremath{r}}
\newcommand{\trapdoorext}{\ensuremath{r_{\mathrm{ext}}}\xspace}
\newcommand{\trapdoorsim}{\ensuremath{r_{\mathrm{sim}}}\xspace}
\renewcommand{\mod}{\ensuremath{\;\mathrm{mod}\;}}
\newcommand{\hsub}{\ensuremath{H^*}\xspace}
\num\newcommand{\enc1}[1]{\ensuremath{\left[#1\right]_1}\xspace}
\newcommand{\enci}[1]{\ensuremath{\left[#1\right]_i}\xspace}
\num\newcommand{\enc2}[1]{\ensuremath{\left[#1\right]_2}\xspace}
\newcommand{\gen}{\ensuremath{\mathsf{gen}}\xspace}
\newcommand{\prv}{\ensuremath{\mathsf{\mathbf{P}}}\xspace}
\newcommand{\prvpoly}{\ensuremath{\mathrm{P_{\mathsf{poly}}}}\xspace}
\newcommand{\prvpc}{\ensuremath{\mathrm{P_{\mathsf{PC}}}}\xspace}
\newcommand{\verpoly}{\ensuremath{\mathrm{V_{\mathsf{poly}}}}\xspace}
\newcommand{\verpc}{\ensuremath{\mathrm{V_{\mathsf{PC}}}}\xspace}
\newcommand{\ideal}{\ensuremath{\mathcal{I}}\xspace}
\newcommand{\prf}{\ensuremath{\pi}\xspace}
\newcommand{\simprv}{\ensuremath{\mathrm{P^{sim}}}\xspace}

%\newcommand{\enc}[1]{\ensuremath{\left[#1\right]}\xspace}
%\num\newcommand{\G0}{\ensuremath{\mathbf{G}}\xspace}
\newcommand{\GG}{\ensuremath{\mathbf{G^*}}\xspace}  % would have liked to call this G01 but problem with name
\num\newcommand{\g0}{\ensuremath{\mathbf{g}}\xspace}
\newcommand{\inst}{\ensuremath{\phi}\xspace}
\newcommand{\inp}{\ensuremath{\mathsf{x}}\xspace}
\newcommand{\wit}{\ensuremath{\omega}\xspace}
\newcommand{\ver}{\ensuremath{\mathsf{\mathbf{V}}}\xspace}
\newcommand{\per}{\ensuremath{\mathsf{\mathbf{P}}}\xspace}
\newcommand{\sonic}{\ensuremath{\mathsf{Sonic}}\xspace}
\newcommand{\aurora}{\ensuremath{\mathsf{Aurora}}\xspace}
\newcommand{\auroralight}{\ensuremath{\mathsf{Auroralight}}\xspace}
\newcommand{\groth}{\ensuremath{\mathsf{Groth'16}}\xspace}
\newcommand{\kate}{\ensuremath{\mathsf{KZG}}\xspace}
\newcommand{\rel}{\ensuremath{\mathcal{R}}\xspace}
\newcommand{\ext}{\ensuremath{E}\xspace}
\newcommand{\params}{\ensuremath{\mathsf{params}_{\inst}}\xspace}
\newcommand{\protparams}{\ensuremath{\mathsf{params}_{\inst}^\advv}\xspace}
\num\newcommand{\p1}{\ensuremath{P_1}\xspace}
\newcommand{\advv}{\ensuremath{ {\mathcal A}^{\mathbf{*}}}\xspace} % the adversary that uses protocol adversary as black box
\newcommand{\crs}{\ensuremath{\sigma}\xspace}
%\num\newcommand{\crs1}{\ensuremath{\mathrm{\sigma}_1}\xspace}
%\num\newcommand{\crs2}{\ensuremath{\mathrm{\sigma}_2}\xspace}
\newcommand{\set}[1]{\ensuremath{\left\{#1\right\}}\xspace}
\newcommand{\hgen}{\ensuremath{\mathbf{g}}\xspace}
\renewcommand{\sim}{\ensuremath{\mathsf{sim}}\xspace}%the distribution of messages when \advv simulates message of \p1
\newcommand{\real}{\ensuremath{\mathsf{real}}\xspace}%the distribution of messages when \p1 is honest and \adv controls rest of players
 \newcommand{\koevec}[2]{\ensuremath{(1,#1,\ldots,#1^{#2},\alpha,\alpha #1,\ldots,\alpha #1^{#2})}\xspace}
\newcommand{\mida}{\ensuremath{A_{\mathrm{mid}}}\xspace}
\newcommand{\midb}{\ensuremath{B_{\mathrm{mid}}}\xspace}
\newcommand{\midc}{\ensuremath{C_{\mathrm{mid}}}\xspace}
\newcommand{\chal}{\ensuremath{\mathsf{challenge}}\xspace}
\newcommand{\attackparams}{\ensuremath{\mathsf{params^{pin}}}\xspace}
\newcommand{\pk}{\ensuremath{\mathsf{pk}}\xspace}
\newcommand{\attackdist}[2]{\ensuremath{AD_{#1}}\xspace}
\renewcommand{\neg}{\ensuremath{\mathsf{negl}(\lambda)}\xspace}
\newcommand{\ro}{\ensuremath{{\mathcal R}}\xspace}
\newcommand{\elements}[1]{\ensuremath{\mathsf{elements}_{#1}}\xspace}
 \num\newcommand{\elmpowers1}[1]{\ensuremath{\mathrm{\mathsf{e}}^1_{#1}}\xspace}
 \num\newcommand{\elmpowers2}[1]{\ensuremath{\mathrm{\mathsf{e}}^2_{#1}}\xspace}
\newcommand{\elempowrs}[1]{\ensuremath{\mathsf{e}_{#1}}\xspace}
 \newcommand{\secrets}{\ensuremath{\mathsf{secrets}}\xspace}
 \newcommand{\polysofdeg}[1]{\ensuremath{\F_{< #1}[X]}\xspace}
 \newcommand{\bivar}[1]{\ensuremath{\F_{< #1}[X,Y]}\xspace}
 \newcommand{\sig}{\ensuremath{\mathscr{S}}\xspace}
 \newcommand{\prot}{\ensuremath{\mathscr{P}}\xspace}
 \newcommand{\PCscheme}{\ensuremath{\mathscr{S}}\xspace}
 \newcommand{\protprime}{\ensuremath{\mathscr{P^*}}\xspace}
 \newcommand{\sigprv}{\ensuremath{\mathsf{P_{sc}}}\xspace}
 \newcommand{\sigver}{\ensuremath{\mathsf{V_{sc}}}\xspace}
 \newcommand{\sigpoly}{\ensuremath{\mathsf{S_{\sigma}}}\xspace}
 \newcommand{\idpoly}{\ensuremath{\mathsf{S_{ID}}}\xspace}
\newcommand{\idpolyevala}{\ensuremath{\mathsf{\bar{s}_{ID1}}}\xspace}
\newcommand{\sigpolyevala}{\ensuremath{\mathsf{\bar{s}_{\sigma1}}}\xspace}
\newcommand{\sigpolyevalb}{\ensuremath{\mathsf{\bar{s}_{\sigma2}}}\xspace}
\newcommand{\bctv}{\ensuremath{\mathsf{BCTV}}\xspace}
\newcommand{\PI}{\ensuremath{\mathsf{PI}}\xspace}
\newcommand{\PIb}{\ensuremath{\mathsf{PI_B}}\xspace}
\newcommand{\PIc}{\ensuremath{\mathsf{PI_C}}\xspace}
\newcommand{\dl}[1]{\ensuremath{\widehat{#1}}\xspace}
\newcommand{\obgen}{\ensuremath{\mathcal O}\xspace}
\newcommand{\PC}{\ensuremath{\mathscr{P}}\xspace}
\newcommand{\permscheme}{\ensuremath{\sigma_\mathscr{P}}\xspace}

	
%\newcommand{\open}{\mathsf{open}}
\newcommand{\update}{\mathsf{update}}
\newcommand{\Prove}{\mathcal{P}}
\newcommand{\Verify}{\mathcal{V}}
\newcommand{\Extract}{\mathcal{E}}
\newcommand{\Simulate}{\mathcal{S}}
\newcommand{\Unique}{\mathcal{U}}
\newcommand{\Rpoly}{\R{\poly}}
\newcommand{\Ppoly}{\Prove{\poly}}
\newcommand{\Vpoly}{\Verify{\poly}}
\newcommand{\Psnark}{\prv}%{\Prove{\snark}}
\newcommand{\Vsnark}{\ver}%{\Verify{\snark}}
\newcommand{\Rprod}{\R{\grandprod}}
\newcommand{\Pprod}{\Prove{\grandprod}}
\newcommand{\Vprod}{\Verify{\grandprod}}
\newcommand{\Rperm}{\R{\perm}}
\newcommand{\Pperm}{\Prove{\perm}}
\newcommand{\Vperm}{\Verify{\perm}}
\newcommand{\f}{\ensuremath{\F}\xspace}
\newcommand{\fn}{\ensuremath{\F^n}\xspace}
\newcommand{\gn}{\ensuremath{\G^n}\xspace}
\newcommand{\fnn}{\ensuremath{\F^{n+1}}\xspace}
% \newcommand{\zw}[1]{{\textcolor{magenta}{Zac:#1}}}
% \newcommand{\ag}[1]{{\textcolor{blue}{\emph{Ariel:#1}}}}
\newtheorem{lemma}{Lemma}[section]
\newtheorem{thm}{Theorem}
\newtheorem{dfn}[lemma]{Definition}
\newtheorem{remark}[lemma]{Remark}


\newtheorem{claim}[lemma]{Claim}
\newtheorem{corollary}[lemma]{Corollary}

\newcommand{\inprod}[2]{\ensuremath{\left\langle #1,#2 \right \rangle}\xspace}
\newcommand{\sumi}[1]{\sum_{i< #1}}
\newcommand{\prob}{\ensuremath{\mathsf{prob}}\xspace}
\begin{document}
    \maketitle

\begin{abstract}
Based on the formal verification of the Zcash Ironwood protocol, we present a security proof of the inner product argument (IPA) in the algebraic group model that is arguably simpler than previous proofs.

\end{abstract}

\section{Introduction}
The IPA \cite{bootle,bulletproofs} is important.
For example, for Zcash Orchard and Ironwood.
The correctness of the IPA is tricky. A subtle mix of ``split-and-fold'' sumcheck like protocol and the discrete log hardness. 
Ironwood FV has nice proof leveraging the Algebraic Group Model (AGM).
Purpose of this note is to have as simple as possible description of IPA protocol and its knowledge soundness in human readable form.
\section{Preliminaries}
\paragraph{Terminology}
We assume our field \F and group \G are both of the same prime size $r$. $\G$ is written additively. Hence, we can write $aG$ for $a\in \F,G \in \G$.
It will be convenient to denote inner product of vectors as regular multiplication. Namely,
for $a,b\in \fn$ we denote $ab=a\cdot b \defeq \sum_{i=1}^n a_i b_i$.
Similarly, for $a\in \fn, G\in \gn$ we denote  $aG=a\cdot b \defeq \sum_{i=1}^n a_i G_i$.
For $u\in \f, a \in \fn$ or $a\in \gn$ we denote $ua\defeq (ua_1,\ldots, ua_n)\in \fn$
\paragraph{The discrete log assumption}
We state the discrete log assumption in a multi-generator form that will be convenient.
It is standard that this follows for the regular discrete-log assumption
\begin{dfn}\label{dfn:DLA}[DLA assumption for \G]
For any efficient \adv the probability of winning the following game is \negl.
\begin{enumerate}
 \item Uniformly chosen $G_1,\ldots,G_n\in \G$ are sent to \adv.
 \item \adv outputs a vector $f\in \F^n$.
 \item \adv wins if and only if  $f\neq 0^n$, and   $\inprod{f}{(G_1,\ldots,G_n)}=0$.
\end{enumerate}
\end{dfn}
\section{AGM stuff}
Need to define discrepancy of expressions.

\adv outputs pair $(C,a,v)$ with $\sumi{a_i G_i + vP=C}$
Given an expression $\alpha C +\beta D$, $C,D$ being group elems for \adv or expressions in them
we have 
\newcommand{\vecs}{\ensuremath{\F^n}\xspace}
\section{Inner product opening scheme }
We give a definit
\begin{dfn}\label{dfn:PCscheme}
A linear commitment scheme consists of 
\begin{itemize}
 \item $\gen(n)$ -  a randomized algorithm that outputs a common reference string \srs.
 \item $\com(f,\srs)$ - that given a vector $a\in \vecs$ returns a commitment \cm to $f$.
 \item A public coin protocol  \open between parties \prv and \ver. \prv and \ver are both given $\cm\in \G,b\in \vecs,v\in \F$. \prv is given $a\in \vecs$. 
 \end{itemize}
such that
\begin{itemize}
 \item \textbf{Completeness:} Fix $\cm\in \G, b\in \fn, z\in \F$ and $a$.
 Suppose that  $\cm = \com(a,\srs)$ and $a\cdot b=z$.  Then if \open is run correctly with inputs $(\cm,b,z;a)$
\ver outputs \acc with probability one.
 \item \textbf{Knowledge soundness in the algebraic group model:} There exists an efficient \ext such that for any algebraic adversary \adv the probability of \adv winning the following game is \negl over the randomness of \adv and \gen.
 \begin{enumerate}
  \item Given \srs, \adv outputs $\cm$.
  \item\label{step:ksrndone} \ext, given access to the messages of \adv during the previous step, outputs $a\in \fn$.
  \item\label{step:ksrndtwo} \adv outputs $ b\in \fn$, $z \in \F$.
  \item \adv takes the part of \prv in the protocol \open with joint inputs $(\cm,b,z)$.
  \item $\adv$ wins if 
  \begin{itemize} 
   \item \verpc outputs \acc at the end of the protocol.
   \item $ab\neq z$ or $aG\neq \cm$.
  \end{itemize}

 \end{enumerate}

\end{itemize}
\end{dfn}


\newcommand{\disc}[1]{\ensuremath{\mathsf{disc}(#1)}\xspace}
\newcommand{\alo}{\ensuremath{a_{i,\mathsf{lo}}}\xspace}
\newcommand{\anext}{\ensuremath{a_{i+1}}\xspace}
\newcommand{\bnext}{\ensuremath{b_{i+1}}\xspace}
\newcommand{\gnext}{\ensuremath{G_{i+1}}\xspace}
\newcommand{\cnext}{\ensuremath{C_{i+1}}\xspace}
\newcommand{\ahi}{\ensuremath{a_{i,\mathsf{hi}}}\xspace}
\newcommand{\glo}{\ensuremath{G_{i,\mathsf{lo}}}\xspace}
\newcommand{\ghi}{\ensuremath{G_{i,\mathsf{hi}}}\xspace}
\newcommand{\blo}{\ensuremath{b_{i,\mathsf{lo}}}\xspace}
\newcommand{\bhi}{\ensuremath{b_{i,\mathsf{hi}}}\xspace}
\newcommand{\sumk}[1]{\ensuremath{\sum_{i=1}^k #1}\xspace}
\newcommand{\expr}{\sum_{i=1}^m \gamma_i C_i}
\newcommand{\todo}[1]{TODO:#1}
\section{The IPA protocol}
\paragraph{Notation}
Let $k=\log n$.
$G\in \fn$ is a set of $n$ generators. $G'\defeq (G,P)\in \fnn$.

\noindent
\underline{$\gen(n)$:}
Choose uniform independent elements $G\in \gn$ and $P\in \G$.
Output $\srs=(G,P)$.\\

\noindent
\underline{$\com(a\in \fn,\srs=(G,P))$:}
Compute and output $\cm\defeq aG$. \\ \\ 
\noindent
\underline{$\open(\cm,v,b;a)$}
\begin{enumerate}
 \item \ver sends random $\alpha\in \F$. 
\item Denote $a_0\defeq a, b_0\defeq b, C_0\defeq \cm + v \alpha P$.
 \item For $i=0,\ldots,k-1$
 \begin{enumerate}
  \item \prv computes and sends the terms 
  \[L_i\defeq \ahi(\glo+\blo P), R_i\defeq \alo(\ghi +\bhi P).\]
  \item \ver sends a random $u_i\in \F$.
  
  \item \prv and \ver compute the terms for the next round:
  \begin{itemize}
 \item $\gnext \defeq \glo + u_i \ghi$,
 \item $\bnext \defeq \blo + u_i\bhi$,
 \item $\cnext \defeq (1/u_i) L_i + C_i + u_i R_i$.
  \end{itemize}
 \item \prv also computes $\anext \defeq \alo + (1/u_i) \ahi$.
 \end{enumerate}
% \item Let $a_k,b_k,G_k,C_k$ be the terms  from the last round.
\item \prv sends $a_k$ to \ver.
\item \ver outputs \acc if and only if
\[C_k= a_k G_k + a_k b_k \alpha P\]
\end{enumerate}

\section{Analyzing knowledge soundness in the AGM}
\paragraph{Discrepancy of expressions}
Let $C$ be any group element output by \adv during the KS game.
As \adv is algebraic $C$ is accompanied by an explanation $(a,v')$ with $a\in \fn,v \in \F$, such that $C=aG+vP$.


We define the discrepancy 
\[\disc{C}\defeq ab-v/\alpha.\]
Elements of discrepancy zero are those where the coefficient $v/\alpha$ of $P'$ equals the inner product $ab$ where $a$ is the coefficient vector for $G$.
In particular, $\disc{C_0}=0$ when \prv is honest.

By linearity, we can also define $\disc{C}$ for any \emph{expression} $\sum_{i=1}^m \gamma_i C_i$ computed by \ver during \open.
% for some integer $m$, $\gamma_i \in \F$ and elements $C_i$ output by 
% \adv during the game.
Namely, we define
\[\disc{\sum_{i=1}^m \gamma_i C_i}\defeq \sum_{i=1}^m \gamma_i \disc{C_i}.\]



The crucial point for us is that a expression evaluating to zero with non-zero discrepancy gives us a discrete log relation:
Suppose  $\expr=0$. For each $i\in [m]$ we can write 


\begin{claim}\label{clm:dlafromdisc}
Given an explained element $C\in \G$ with $C=0$ and $\disc{C}\neq 0$
 we get a DL solution $a \in \fn$.
\end{claim}
\begin{proof}
Let $a'=(a,v)\in \fnn$ be the explanation of $C$. $\disc{C}\neq 0$ means $ab -v/\alpha \neq 0$.
In particular, this implies $a'\neq 0$.
Since $a'G'=0$ we have that $a'$ is a DL solution for $G'$.
\end{proof}


\subsection{The main argument}
We first define the extractor \ext:
Let $(a,v)$ be the explanation \adv sends along with \cm in the first step of the game.
That is $\cm = aG+vP$.
\ext simply outputs $a$.


Let $E$ be the event that \adv wins the knowledge-soundness game.
Namely, the event that
-
% that is \ver accepts while $ab\neq v$, where
\begin{itemize}
   \item \ver outputs \acc at the end of the protocol.
   \item $ab\neq z$ or $aG\neq \cm$.
\end{itemize}
We need to show $\prob(E)=\negl$.
The basic idea is:
 Let $Z$ be the final term \ver checks is zero; i.e.
 \[Z\defeq C_k - a_k G_k - a_k b_k \alpha P.\]
We show that except for a \negl subevent of $E$, when we are in $E$, \disc{Z} is non-zero.
Which means according to Claim \ref{clm:dlafromdisc} we can output a DLA relation. We proceed with the details.

We have $C_0 =\cm+z\alpha P$.
So $C_0$'s explanation is $(a,v+\alpha z)$
Let $E_0\subset E$ be the subevent of $E$ where $\disc{C_0}=0$. We claim that $\prob(E_0)=\negl$:
% We look at two cases
% \begin{itemize}
%  \item $ab\neq z$:
 \[\disc{C_0} = ab  - v/\alpha - z\]
If $v\neq 0$ the probability over $\alpha$ that this is zero is at most $1/r=\negl$.
If $v=0$ it means $aG=\cm$ so we're in the event that $ab\neq z$ so $\disc{C_0}\neq 0$.


% \end{itemize}
Now assume we are in $E\setminus E_0$ so that $\disc{C_0}\neq 0$.
For $i\in [k]$ let $E_i$ be the event that $\disc{C_i}= 0$ while $\disc{C_j}\neq 0$ for $j<i$.
From linearity, we have
\[\disc{C_i}= (1/u)\disc{L_i}+ \disc{C_{i-1}} + u\disc{R_i}.\]
Thus, when \disc{C_{i-1}} is non-zero the probability that $\disc{C_i}=0$ is at most $2/r$.
So for all $i\in [k]$, $\prob(E_i)=\negl$.
When we are outside the event $E_1\cup\ldots\cup E_k$  (but in $E\setminus E_0$) then 
\[\disc{C_k}\neq 0.\]

Let $U\in \fn$ be the tensor vector 
\[U\defeq (1,u_1)\otimes \ldots\otimes (1,u_k)\]
The important observation is that both $b_k$ and $G_k$ are inner products with $U$:
\[b_k=Ub,G_k=UG.\]
Thus
\[\disc{a_kG_k - a_k b_k\alpha P} = a_k\cdot \disc{UG-(Ub)\alpha P}=0.\]
Let $Z$ be the final term \ver checks is zero; i.e.
\[Z\defeq C_k - a_k G_k - a_k b_k \alpha P.\]
So we have 

So
\[\disc{Z}=\disc{C_k}- \disc{a_k G_k - a_k b_k \alpha P} = \disc{C_k}\neq 0\]

And so in the event $E\setminus (E_1\cup\ldots\cup  E_k)$ we can output a DL solution for $G'$ according to Claim \ref{clm:dlafromdisc}.
Hence, according to DLA as in Definition \ref{dfn:DLA} the probability of $E\setminus (E_1\cup\ldots\cup E_k)$ is \negl.
And so $\prob(E)=\negl$.
% \end{proof}

% An \emph{expression} is either an element output by \adv during the knowledge soundness game
 \section*{Acknowledgements}   We thank ZCG for supporting this work
 \bibliographystyle{alpha}
\bibliography{references}
    
\end{document}


